\exercisesection{\S~7}

All the g-modules considered (except those in Exerc. 14 and 15) are assumed to be finite dimensional.

\begin{enumerate}
\def\labelenumi{\arabic{enumi})}
\tightlist
\item
  Let \(\omega \in \mathrm{P}_{++}\) ; denote by \(\mathrm{S}(\omega)\) the set of weights of \(\mathrm{E}(\omega)\) , in other words the smallest R-saturated subset of P containing \(\omega\) (Prop. 5). If \(\lambda \in \mathrm{S}(\omega)\) , we have \(\lambda \equiv \omega\) (mod. Q). Conversely, let \(\lambda \in \mathrm{P}\) be such that \(\lambda \equiv \omega\) (mod. Q); prove the equivalence of the following properties:
\end{enumerate}

\begin{enumerate}
\def\labelenumi{(\roman{enumi})}
\item
  \(\lambda \in S(\omega);\)
\item
  \(\omega - w\lambda \in Q_{+}\) for all \(w \in W\) ;
\item
  \(\lambda\) belongs to the convex hull of W. \(\omega\) in \(\mathfrak{h}_{\mathbf{R}}^{*}\) .
\end{enumerate}

(To prove that (iii) \(\Longrightarrow\) (ii), remark that \(\omega - w\omega\) is a linear combination of elements of \(\mathrm{R}_{+}\) with coefficients \(\geq 0\) ; deduce, by convexity, that \(\omega - w\lambda\) has the same property; since \(\omega - w\lambda\) belongs to Q, this implies that \(\omega - w\lambda \in \mathrm{Q}_{+}\) . To prove that (ii) \(\Longrightarrow\) (i), choose \(w\) such that \(w\lambda \in \mathrm{P}_{++}\) , and apply Cor. 2 of Prop. 3. The implication (i) \(\Longrightarrow\) (iii) is immediate.)

\begin{enumerate}
\def\labelenumi{\arabic{enumi})}
\setcounter{enumi}{1}
\tightlist
\item
  Let \((\mathrm{R}_{i})_{i\in\mathrm{I}}\) be the family of irreducible components of R, and \(g=\prod_{i\in I}g_{i}\) the corresponding decomposition of g into a product of simple algebras. Identify P with the product of the \(\mathrm{P}(\mathrm{R}_{i})\) , and give each \(\mathrm{P}(\mathrm{R}_{i})\) the order relation defined by the basis \(B_{i}=B\cap R_{i}\) .
\end{enumerate}

\begin{enumerate}
\def\labelenumi{\alph{enumi})}
\item
  Let \(\omega = (\omega_{i})_{i\in \mathrm{I}}\) be an element of \(\mathrm{P}_{++} = \prod_{i\in \mathrm{I}}\mathrm{P}_{++}(\mathrm{R}_i)\) . Show that the simple \(\mathfrak{g}\) -module \(\operatorname {E}(\omega)\) is isomorphic to the tensor product of the simple \(\mathfrak{g}_i\) -modules \(\operatorname {E}(\omega_i)\) .
\item
  Let \(\mathcal{M}\) (resp. \(M_{i}\) ) be the set of elements of \(P_{++}\) (resp. of \(P_{++}(R_{i})\) ) having the equivalent properties (i), (ii), (iii), (iv) of Prop. 6 and 7. Show that \(M = \prod_{i \in I} M_{i}\) , in other words that \(\omega \in M\) if and only if, for all \(i \in I\) , \(\omega_{i}\) is either zero or a minuscule weight of \(R_{i}\) . Deduce that M is a system of representatives in P of the elements of P/Q.
\item
  Let E be a simple g-module, and X its set of weights. Show that X contains a unique element of M, and that the multiplicity of this element is equal to the upper bound of the multiplicities of the elements of X.
\end{enumerate}

\begin{enumerate}
\def\labelenumi{\arabic{enumi})}
\setcounter{enumi}{2}
\tightlist
\item
  \begin{enumerate}
  \def\labelenumii{\alph{enumii})}
  \tightlist
  \item
    Let E be a g-module. Show the equivalence of the conditions:
  \end{enumerate}
\end{enumerate}

\begin{enumerate}
\def\labelenumi{(\roman{enumi})}
\item
  The rank of the semi-direct product of \(\mathfrak{g}\) by E is strictly larger than that of \(\mathfrak{g}\) .
\item
  0 is a weight of E.
\item
  There exists a weight of \(\mathrm{E}\) that is radical (i.e.~belongs to \(\mathrm{Q}\) ).
\end{enumerate}

\begin{enumerate}
\def\labelenumi{\alph{enumi})}
\setcounter{enumi}{1}
\tightlist
\item
  Assume that \(\mathrm{E}\) is simple. Show that (i), (ii), (iii) are equivalent to
\end{enumerate}

\begin{enumerate}
\def\labelenumi{(\roman{enumi})}
\setcounter{enumi}{3}
\tightlist
\item
  The highest weight of E is radical.
\end{enumerate}

If these conditions are satisfied, there exists no non-zero invariant alternating bilinear form. (Use Prop. 12 and Prop. 1 (ii) of §6.)

\begin{enumerate}
\def\labelenumi{\arabic{enumi})}
\setcounter{enumi}{3}
\item
  Let \(k'\) be an extension of \(k\) , and \(\mathfrak{g}' = \mathfrak{g}_{(k')}\) . Show that every \(\mathfrak{g}'\) -module arises, by extension of scalars, from a \(\mathfrak{g}\) -module that is unique up to isomorphism.
\item
  Let \(\mathbf{E}\) be a \(\mathfrak{g}\) -module. Show the equivalence of the conditions:
\end{enumerate}

\begin{enumerate}
\def\labelenumi{(\roman{enumi})}
\item
  E is faithful (i.e.~the canonical map from g to gl(E) is injective).
\item
  Every root of g is the difference between two weights of E.
\end{enumerate}

¶6) Let \(\varphi\) be an involutive automorphism of \(\mathfrak{g}\) whose restriction to \(\mathfrak{h}\) is -Id. Show that, if E is a \(\mathfrak{g}\) -module, there exists a non-degenerate symmetric bilinear form \(\varPsi\) on E such that

\[
\varPsi (x. a, b) + \varPsi (a, \varphi (x). b) = 0 \quad \text {   for   } x \in \mathfrak {g}, a, b \in \mathrm{E}.
\]

(Reduce to the case in which E is simple. Show that the transform of E by \(\varphi\) is isomorphic to the dual \(E^{*}\) of E, and deduce the existence of a non-degenerate bilinear form \(\Psi\) satisfying the conditions above. Show next that, if e is a primitive vector in E, then \(\Psi(e,e)\neq0\) . Deduce that \(\Psi\) is symmetric.)

\begin{enumerate}
\def\labelenumi{\arabic{enumi})}
\setcounter{enumi}{6}
\tightlist
\item
  If \(\lambda \in \mathrm{P}_{++}\) , denote by \(\rho_{\lambda}:\mathrm{U}(\mathfrak{g})\to \mathrm{End}(\mathrm{E}(\lambda))\) the representation defined by the simple module \(\mathrm{E}(\lambda)\) . We have \(\mathrm{Im}(\rho_{\lambda}) = \mathrm{End}(\mathrm{E}(\lambda));\) put \(\mathfrak{m}_{\lambda} = \mathrm{Ker}(\rho_{\lambda})\) . \(a)\) Show that the \(\mathfrak{m}_{\lambda}\) are pairwise distinct, and that they are the only two-sided ideals \(\mathfrak{m}\) of \(\mathrm{U}(\mathfrak{g})\) such that \(\mathrm{U}(\mathfrak{g}) / \mathfrak{m}\) is a finite dimensional simple \(k\) -algebra.
\end{enumerate}

\begin{enumerate}
\def\labelenumi{\alph{enumi})}
\setcounter{enumi}{1}
\item
  If I is a finite subset of \(P_{++}\) , put \(m_{I} = \bigcap_{\lambda \in I} m_{\lambda}\) . Show that the canonical map \(U(\mathfrak{g}) / \mathfrak{m}_{I} \to \prod_{\lambda \in I} U(\mathfrak{g}) / \mathfrak{m}_{\lambda}\) is an isomorphism, and that every two-sided ideal of \(U(\mathfrak{g})\) of finite codimension is equal to exactly one of the \(m_{I}\) .
\item
  Show that the principal anti-automorphism of \(\mathrm{U}(\mathfrak{g})\) transforms \(\mathfrak{m}_{\lambda}\) to \(\mathfrak{m}_{\lambda^*}\) , where \(\lambda^{*} = -w_{0}\lambda\) (cf.~Prop. 11).
\end{enumerate}

¶8) Let g be a semi-simple Lie algebra; exceptionally, we do not assume in this exercise that g is split. Let \(\bar{k}\) be an algebraic closure of k and let

\[
\pi : \operatorname{Gal} (\bar {k} / k) \to \operatorname{Aut} (\bar {\mathrm{R}}, \bar {\mathrm{B}})
\]

be the homomorphism defined in §5, Exerc. 8. Let \(\operatorname{Gal}(\bar{k}/k)\) operate, via \(\pi\) , on the set \(\bar{P}_{++}\) of dominant weights of \(\bar{R}\) relative to \(\bar{B}\) ; let \(\Omega\) be a system of representatives of the elements of the quotient \(\bar{P}_{++}/\operatorname{Gal}(\bar{k}/k)\) .

\begin{enumerate}
\def\labelenumi{\alph{enumi})}
\item
  Put \(\bar{g} = \bar{k} \otimes_{k} g\) . If I is a finite subset of \(\bar{P}_{++}\) stable under \(\operatorname{Gal}(\bar{k}/k)\) , the two-sided ideal \(\bar{m}_{I}\) of \(\operatorname{U}(\bar{g})\) associated to I (cf.~Exerc. 7) is of the form \(\bar{k} \otimes_{k} m_{I}\) , where \(m_{I}\) is a two-sided ideal of \(\operatorname{U}(g)\) . Show that every two-sided ideal of \(\operatorname{U}(g)\) of finite codimension is obtained in this way exactly once.
\item
  Let \(\omega \in \Omega\) , \(\mathrm{I}(\omega)\) its orbit under \(\operatorname{Gal}(\bar{k}/k)\) , and \(\mathrm{G}_{\omega}\) the stabilizer of \(\omega\) ; let \(k_{\omega}\) be the sub-extension of \(\bar{k}\) corresponding to \(\mathrm{G}_{\omega}\) by Galois theory. Show that \(\mathrm{U}(\mathfrak{g})/\mathfrak{m}_{\mathrm{I}(\omega)}\) is a simple algebra whose centre is isomorphic to \(k_{\omega}\) . Every two-sided ideal \(\mathfrak{m}\) of \(\mathrm{U}(\mathfrak{g})\) such that \(\mathrm{U}(\mathfrak{g})/\mathfrak{m}\) is a finite dimensional simple algebra is equal to exactly one of the \(\mathfrak{m}_{\mathrm{I}(\omega)}\) .
\item
  The group \(\operatorname{Gal}(\bar{k}/k)\) operates, via \(\pi\) , on the ring \(\operatorname{R}(\mathfrak{g})\) ; denote by \(\operatorname{R}(\mathfrak{g})^{\mathrm{inv}}\) the subring of \(\operatorname{R}(\mathfrak{g})\) consisting of the elements invariant under \(\operatorname{Gal}(\bar{k}/k)\) . Show that the map \([E] \to [\bar{k} \otimes_{k} E]\) extends to an injective homomorphism from \(\operatorname{R}(\mathfrak{g})\) to \(\operatorname{R}(\mathfrak{g})^{\mathrm{inv}}\) whose cokernel is a torsion group; this is an isomorphism if and only if, for all \(\omega \in \Omega\) , \(\operatorname{U}(\mathfrak{g})/\mathfrak{m}_{\operatorname{I}(\omega)}\) is an algebra of matrices over \(k_{\omega}\) ; show that this is the case when g has a Borel subalgebra. \(^{14}\)
\end{enumerate}

\begin{enumerate}
\def\labelenumi{\arabic{enumi})}
\setcounter{enumi}{8}
\tightlist
\item
  Let \(\mathfrak{a}_1\) and \(\mathfrak{a}_2\) be Lie algebras. Show that there exists a unique homomorphism
\end{enumerate}

\[
f: \mathrm{R} (\mathfrak {a} _ {1}) \otimes_ {\mathbf {Z}} \mathrm{R} (\mathfrak {a} _ {2}) \to \mathrm{R} (\mathfrak {a} _ {1} \times \mathfrak {a} _ {2})
\]

such that \(f([\mathrm{E}_1] \otimes [\mathrm{E}_2]) = [\mathrm{E}_1 \otimes \mathrm{E}_2]\) if \(\mathrm{E}_i\) is an \(\mathfrak{a}_i\) -module ( \(i = 1, 2\) ). Show that \(f\) is injective, and that it is bijective if \(\mathfrak{a}_1\) and \(\mathfrak{a}_2\) are splittable semi-simple.

\begin{enumerate}
\def\labelenumi{\arabic{enumi})}
\setcounter{enumi}{9}
\tightlist
\item
  Let \(\Gamma\) be a subgroup of \(P\) containing \(Q\) . Such a subgroup is stable under \(W\) .
\end{enumerate}

\begin{enumerate}
\def\labelenumi{\alph{enumi})}
\item
  Show that, if \(\lambda \in \mathrm{P}_{++} \cap \Gamma\) , every weight of \(\operatorname{E}(\lambda)\) belongs to \(\Gamma\) .
\item
  Let \(\mathrm{R}_{\Gamma}(\mathfrak{g})\) be the subgroup of \(\mathrm{R}(\mathfrak{g})\) with basis the \([\lambda]\) , with \(\lambda \in \mathrm{P}_{++} \cap \Gamma\) . If \(\mathrm{E}\) is a \(\mathfrak{g}\) -module, \([\mathrm{E}] \in \mathrm{R}_{\Gamma}(\mathfrak{g})\) if and only if the weights of \(\mathrm{E}\) belong to \(\Gamma\) . Deduce that \(\mathrm{R}_{\Gamma}(\mathfrak{g})\) is a subring of \(\mathrm{R}(\mathfrak{g})\) .
\item
  Show that the homomorphism \(\operatorname{ch} : \mathrm{R}_{\Gamma}(\mathfrak{g}) \to \mathbf{Z}[\Gamma]\) is an isomorphism from \(\mathrm{R}_{\Gamma}(\mathfrak{g})\) to the subring of \(\mathbf{Z}[\Gamma]\) consisting of the elements invariant under W. (Use Th. 2 (ii).)
\end{enumerate}

\(d\) ) Describe \(\mathrm{R}(\mathfrak{g})\) and \(\mathrm{R}_{\Gamma}(\mathfrak{g})\) explicitly when \(\mathfrak{g} = \mathfrak{sl}(2,k)\) and \(\Gamma = \mathrm{Q}\) .

¶11) The notations are those of no. 7. For any integer \(m \geq 1\) , denote by \(\Psi^{m}\) the endomorphism of \(\mathbf{Z}[\Delta]\) that takes \(e^{\lambda}\) to \(e^{m\lambda}\) . We have \(\Psi^{1} = \mathrm{Id}\) and \(\Psi^{m} \circ \Psi^{n} = \Psi^{mn}\) .

Let E be a finite dimensional \(\Delta\) -graded vector space. For all \(n \geq 0\) , denote by \(a_{n}E\) (resp. \(s_{n}E\) ) the nth exterior (resp. symmetric) power of E, equipped with its natural grading.

\begin{enumerate}
\def\labelenumi{\alph{enumi})}
\tightlist
\item
  Show that
\end{enumerate}

\[
n \operatorname{ch} (s _ {n} \mathrm{E}) = \sum_ {m = 1} ^ {n} \Psi^ {m} (\operatorname{ch} (\mathrm{E})) \operatorname{ch} (s _ {n - m} \mathrm{E})
\]

and

\[
n \operatorname{ch} (a _ {n} \mathrm{E}) = \sum_ {m = 1} ^ {n} (- 1) ^ {m - 1} \varPsi^ {m} (\operatorname{ch} (\mathrm{E})) \operatorname{ch} (a _ {n - m} \mathrm{E}).
\]

Deduce that \(\operatorname{ch}(s_{n}\operatorname{E})\) and \(\operatorname{ch}(a_{n}\operatorname{E})\) can be expressed as polynomials, with rational coefficients, in the \(\varPsi^{m}(\operatorname{ch}(\operatorname{E}))\) , \(1 \leq m \leq n\) . For example:

\[
\mathrm{ch} (s _ {2} \mathrm{E}) = \frac {1}{2} \mathrm{ch} (\mathrm{E}) ^ {2} + \frac {1}{2} \varPsi^ {2} (\mathrm{ch} (\mathrm{E}))
\]

\[
\mathrm{ch} (a _ {2} \mathrm{E}) = \frac {1}{2} \mathrm{ch} (\mathrm{E}) ^ {2} - \frac {1}{2} \varPsi^ {2} (\mathrm{ch} (\mathrm{E})).
\]

\begin{enumerate}
\def\labelenumi{\alph{enumi})}
\setcounter{enumi}{1}
\tightlist
\item
  Prove the following identities (in the algebra of formal power series in a variable T, with coefficients in \(Q[\Delta]\) ):
\end{enumerate}

\[
\sum_ {n = 0} ^ {\infty} \mathrm{ch} (s _ {n} \mathrm{E}) \mathrm{T} ^ {n} = \exp \left\{\sum_ {m = 1} ^ {\infty} \varPsi^ {m} (\mathrm{ch(E)}) \mathrm{T} ^ {m} / m \right\}
\]

and

\[
\sum_ {n = 0} ^ {\infty} \mathrm{ch} (a _ {n} \mathrm{E}) \mathrm{T} ^ {n} = \exp \left\{\sum_ {m = 1} ^ {\infty} (- 1) ^ {m - 1} \varPsi^ {m} (\mathrm{ch(E)}) \mathrm{T} ^ {m} / m \right\}.
\]

\begin{enumerate}
\def\labelenumi{\alph{enumi})}
\setcounter{enumi}{2}
\tightlist
\item
  Assume that \(\Delta\) is a group, so that \(\varPsi^{m}\) can be defined for all \(m\in \mathbf{Z}\) . Show that, if \(\mathrm{E}^*\) is the graded dual of \(\mathrm{E}\) ,
\end{enumerate}

\[
\operatorname{ch} \left(\mathrm{E} ^ {*}\right) = \Psi^ {- 1} (\operatorname{ch} (\mathrm{E})).
\]

\begin{enumerate}
\def\labelenumi{\alph{enumi})}
\setcounter{enumi}{3}
\tightlist
\item
  Identify \(\mathrm{R}(\mathfrak{g})\) , by means of ch, with a subring of Z[P]. Show that \(\mathrm{R}(\mathfrak{g})\) is stable under the \(\varPsi^{m}\) , \(m \in Z\) , and that so is the subring \(\mathrm{R}_{\varGamma}(\mathfrak{g})\) defined in the preceding exercise.
\end{enumerate}

\begin{enumerate}
\def\labelenumi{\arabic{enumi})}
\setcounter{enumi}{11}
\tightlist
\item
  Let \(\lambda \in \mathrm{P}_{++}\) and let \(\mathcal{X}_{\lambda}\) be the set of weights of \(\operatorname{E}(\lambda)\) . Show that
\end{enumerate}

\[
\mathcal {X} _ {\lambda} \subset \lambda - \mathrm{P} _ {+ +}
\]

does not hold in general. (Consider, for example, the adjoint representation of \(\mathfrak{sl}(3,k)\) .)

\begin{enumerate}
\def\labelenumi{\arabic{enumi})}
\setcounter{enumi}{12}
\tightlist
\item
  Let \(\lambda = \sum_{\alpha \in B} a_{\alpha} \alpha\) be an element of \(P_{++}\) . For \(n = 0, 1, \ldots\) , denote by \(\mathcal{X}_n\) the set of weights \(\mu\) of \(E(\lambda)\) such that \(\lambda - \mu\) is the sum of \(n\) elements of \(B\) . Let \(s_n\) be the sum of the multiplicities of the elements of \(\mathcal{X}_n\) (as weights of \(E(\lambda)\) ). Let \(T = 2 \sum_{\alpha \in B} a_{\alpha}\) . Show that:
\end{enumerate}

\begin{enumerate}
\def\labelenumi{\alph{enumi})}
\item
  T is an integer \(\geq 0\) .
\item
  \(s_{n}=0\) for n > T, and \(s_{T-n}=s_{n}\) .
\item
  If r is the integer part of T/2, then \(s_{1} \leq s_{2} \leq \cdots \leq s_{r+1}\) .
\end{enumerate}

¶14) Let \(\lambda \in \mathrm{P}_{++}\) , \(\mathrm{F}_{\lambda}\) be the largest proper submodule of \(Z(\lambda)\) and \(v\) a primitive element of \(Z(\lambda)\) of weight \(\lambda\) , cf.~§6, no. 3. Show that

\[
\mathrm{F} _ {\lambda} = \sum_ {\alpha \in \mathrm{B}} \mathrm{U} (\mathfrak {g}) X _ {- \alpha} ^ {\lambda (H _ {\alpha}) + 1} v = \sum_ {\alpha \in \mathrm{B}} \mathrm{U} (\mathfrak {n} _ {-}) X _ {- \alpha} ^ {\lambda (H _ {\alpha}) + 1} v.
\]

\begin{enumerate}
\def\labelenumi{\arabic{enumi})}
\setcounter{enumi}{14}
\tightlist
\item
  Let \(\lambda \in \mathfrak{h}^*\) and let \(v\) (resp. \(v'\) ) be a primitive element of \(\mathrm{Z}(\lambda)\) (resp. of \(\mathrm{E}(\lambda)\) ) of weight \(\lambda\) . Let \(\mathrm{I}\) (resp. \(\mathrm{I}'\) ) be the annihilator of \(v\) (resp. \(v'\) ) in \(\mathrm{U}(\mathfrak{g})\) .
\end{enumerate}

\[
a) \mathrm{I} = \mathrm{U} (\mathfrak {g}) \mathfrak {n} _ {+} + \sum_ {h \in \mathfrak {h}} \mathrm{U} (\mathfrak {g}) (h - \lambda (h)).
\]

\begin{enumerate}
\def\labelenumi{\alph{enumi})}
\setcounter{enumi}{1}
\item
  I' is the largest left ideal of U(g) distinct from U(g) and containing I.
\item
  If \(\lambda \in \mathrm{P}_{++}\) , then
\end{enumerate}

\[
\mathrm{I} ^ {\prime} = \mathrm{I} + \sum_ {\alpha \in \mathrm{B}} \mathrm{U} (\mathfrak {g}) X _ {- \alpha} ^ {\lambda (H _ {\alpha}) + 1} = \mathrm{I} + \sum_ {\alpha \in \mathrm{B}} \mathrm{U} (\mathfrak {n} _ {-}) X _ {- \alpha} ^ {\lambda (H _ {\alpha}) + 1}.
\]

(Use the preceding exercise.)

¶16) Let V and V' be two g-modules. Then V' is said to be subordinate to V if there exists a linear map \(f : V \to V'\) such that:

\(\alpha)\) \(f\) is surjective; \(\beta)\) the image under \(f\) of a primitive element of \(\mathbf{V}\) is either 0 or a primitive element of \(\mathrm{V}^{\prime}\) ; \(\gamma)\) \(f\) is an \(\mathfrak{n}_{-}\) -homomorphism.

\begin{enumerate}
\def\labelenumi{\alph{enumi})}
\item
  Let \(f: V \to V'\) satisfy \(\alpha\) ), \(\beta\) ), \(\gamma\) ). Let v be a primitive element of V. Then the image under f of the g-submodule generated by v is the g-submodule generated by \(f(v)\) .
\item
  Let \(f: \mathrm{V} \to \mathrm{V}'\) satisfy \(\alpha), \beta), \gamma)\) . Let \(\mathrm{W}\) be a \(\mathfrak{g}\) -submodule of \(\mathrm{V}\) . Then \(f(\mathrm{W})\) is a \(\mathfrak{g}\) -submodule of \(\mathrm{V}'\) that is subordinate to \(\mathrm{W}\) .
  \item
  Let \(V = E_{1} \oplus \cdots \oplus E_{s}\) and \(V' = E_{1}' \oplus \cdots E_{s'}\) be decompositions of V and \(V'\) into sums of simple modules. Then \(V'\) is subordinate to V if and only if \(s' \leq s\) and there exists \(\sigma \in S_{s}\) such that \(E_{i}'\) is subordinate to \(\mathrm{E}_{\sigma(i)}\) for \(i = 1, \ldots, s'\) .
\item
  If \(V'\) is subordinate to V and if V is simple, then \(V'\) is simple or reduced to 0.
\item
  Assume that V and \(V'\) are simple. Let \(\lambda\) and \(\lambda'\) be their highest weights. Then \(V'\) is subordinate to V if and only if \(\lambda'(H_{\alpha}) \leq \lambda(H_{\alpha})\) for all \(\alpha \in B\) . (For the sufficiency, use the preceding exercise.)
\end{enumerate}

\begin{enumerate}
\def\labelenumi{\arabic{enumi})}
\setcounter{enumi}{16}
\tightlist
\item
  Let \(\lambda, \mu \in \mathrm{P}_{++}\) and \(\alpha \in \mathrm{B}\) be such that \(\lambda(H_{\alpha}) \geq 1\) and \(\mu(H_{\alpha}) \geq 1\) . Let \(\mathrm{F} = \mathrm{E}(\lambda) \otimes \mathrm{E}(\mu)\) .
\end{enumerate}

\begin{enumerate}
\def\labelenumi{\alph{enumi})}
\item
  Show that \(\dim \mathrm{F}^{\lambda +\mu} = 1\) and \(\dim \mathrm{F}^{\lambda +\mu -\alpha} = 2\) .
\item
  Show that \(X_{\alpha}: F^{\lambda+\mu-\alpha} \to F^{\lambda+\mu}\) is surjective, and that the non-zero elements of its kernel are primitive (remark that, if \(\beta \in B\) is distinct from \(\alpha\) , \(\lambda + \mu - \alpha + \beta\) is not a weight of F).
\item
  Deduce that \(\mathrm{E}(\lambda)\otimes\mathrm{E}(\mu)\) contains a unique submodule isomorphic to
\end{enumerate}

\[
\mathrm{E} (\lambda + \mu - \alpha).
\]

\begin{enumerate}
\def\labelenumi{\alph{enumi})}
\setcounter{enumi}{3}
\tightlist
\item
  Show that \(\mathbf{S}^{2}(\mathrm{E}(\lambda))\) (resp. \(\bigwedge^{2}\mathrm{E}(\lambda)\) ) contains a unique submodule isomorphic to \(\mathrm{E}(2\lambda)\) (resp. to \(\mathrm{E}(2\lambda-\alpha)\) ).
\end{enumerate}

¶ 18) Choose a positive non-degenerate symmetric bilinear form \((\cdot|\cdot)\) on \(\mathfrak{h}_{\mathbf{R}}^{*}\) invariant under W. Let \(\lambda \in P_{++}\) .

\begin{enumerate}
\def\labelenumi{\alph{enumi})}
\item
  Let \(\mu\) be a weight of \(\mathrm{E}(\lambda)\) . Write \(\lambda - \mu\) in the form \(\sum_{\alpha \in \mathrm{B}} k_{\alpha} \alpha\) . Let \(\alpha \in \mathrm{B}\) be such that \(k_{\alpha} \neq 0\) . Show that there exists \(\alpha_{1}, \ldots, \alpha_{n} \in \mathrm{B}\) such that \((\lambda | \alpha_{1}) \neq 0\) , \((\alpha_{1} | \alpha_{2}) \neq 0, \ldots, (\alpha_{n-1} | \alpha_{n}) \neq 0\) , \((\alpha_{n} | \alpha) \neq 0\) .
\item
  Let \(v\) be a primitive element of \(\mathrm{E}(\lambda)\) , and let \(\alpha_1, \ldots, \alpha_n \in \mathrm{B}\) satisfy the following conditions:
\end{enumerate}

\begin{enumerate}
\def\labelenumi{(\roman{enumi})}
\item
  \((\alpha_{i}|\alpha_{i+1})\neq0\) for \(i=1,2,\ldots,n-1;\)
\item
  \((\alpha_{i}|\alpha_{j}) = 0\) for \(j > i + 1\) ;
\item
  \(\lambda(H_{\alpha_{1}})\neq0\) and \(\lambda(H_{\alpha_{2}})=\cdots=\lambda(H_{\alpha_{n}})=0.\)
\end{enumerate}

Show that \(X_{-\alpha_n}X_{-\alpha_{n-1}} \ldots X_{-\alpha_1}v \neq 0\) . (Observe that, for \(1 \leq s \leq n\) ,

\[
\lambda - \alpha_ {1} - \dots - \alpha_ {s - 1} + \alpha_ {n}
\]

is not a weight of \(\mathrm{E}(\lambda)\) , and deduce that \(X_{\alpha_s}X_{-\alpha_s}X_{-\alpha_{s-1}}\ldots X_{-\alpha_1}v\neq 0\) by induction on \(s\) .) If \(\sigma \in \mathfrak{S}_n\) and \(\sigma \neq 1\) , then \(X_{-\alpha_{\sigma(n)}}X_{-\alpha_{\sigma(n-1)}}\ldots X_{-\alpha_{\sigma(1)}}v = 0\) . (Let \(r\) be the smallest integer such that \(\sigma(r)\neq r\) . Use \(a\) ) to show that \(\lambda -\alpha_{\sigma(1)} - \ldots -\alpha_{\sigma(r)}\) is not a weight of \(\mathrm{E}(\lambda)\) .)

\begin{enumerate}
\def\labelenumi{\alph{enumi})}
\setcounter{enumi}{2}
\tightlist
\item
  Let \(\lambda' \in \mathrm{P}_{++}\) . A chain joining \(\lambda\) to \(\lambda'\) is a sequence \((\alpha_1, \ldots, \alpha_n)\) of elements of B such that \(n \geq 1\) , \((\lambda|\alpha_1) \neq 0\) , \((\alpha_1|\alpha_2) \neq 0, \ldots, (\alpha_{n-1}|\alpha_n) \neq 0\) , \((\alpha_n|\lambda') \neq 0\) . Such a chain is called minimal if no strict subsequence of \((\alpha_1, \ldots, \alpha_n)\) joins \(\lambda\) to \(\lambda'\) . In that case, we have \((\alpha_i|\alpha_j) = 0\) if \(|i - j| \geq 2\) , \((\lambda|\alpha_i) = 0\) if \(i \geq 2\) and \((\lambda'|\alpha_i) = 0\) if \(i \leq n - 1\) .
\end{enumerate}

\(d\) ) Let \((\alpha_{1},\ldots ,\alpha_{n})\) be a minimal chain joining \(\lambda\) to \(\lambda^{\prime}\) . If \(v^{\prime}\) is a primitive vector in \(\operatorname {E}(\lambda ')\) , put:

\[
v _ {s} = X _ {- \alpha_ {s}} X _ {- \alpha_ {s - 1}} \ldots X _ {- \alpha_ {1}} v (s = 0, 1, \ldots , n)
\]

\[
v _ {s} ^ {\prime} = X _ {- \alpha_ {s + 1}} X _ {- \alpha_ {s + 2}} \dots X _ {- \alpha_ {n}} v ^ {\prime} (s = 0, 1, \ldots , n)
\]

\(a_0 = (\lambda |\alpha_1), a_n = (-1)^n (\lambda '|\alpha_n), a_s = (-1)^{s + 1}(\alpha_s|\alpha_{s + 1}), 1 \leq s \leq n - 1.\) Show that

\[
\sum_ {s = 0} ^ {n} a _ {s} v _ {s} \otimes v _ {s} ^ {\prime}
\]

is a primitive element of \(\mathrm{E}(\lambda)\otimes\mathrm{E}(\lambda')\) of weight \(\lambda+\lambda'-\alpha_{1}-\cdots-\alpha_{n}\) , and that it is unique up to homothety.

(Use b) and c) to show that every element of \(\mathrm{E}(\lambda) \otimes \mathrm{E}(\lambda')\) of weight

\[
\lambda + \lambda^ {\prime} - \alpha_ {1} - \dots - \alpha_ {n}
\]

is a linear combination of the \(v_{s} \otimes v_{s}^{\prime}\) . Then write down the condition for such a linear combination to be a primitive vector.)

Deduce that \(\mathrm{E}(\lambda) \otimes \mathrm{E}(\lambda')\) contains a unique g-submodule isomorphic to

\[
\operatorname{E} \left(\lambda + \lambda^ {\prime} - \alpha_ {1} - \dots - \alpha_ {n}\right).
\]

(When \(n = 1\) we recover Exerc. 17.)

\begin{enumerate}
\def\labelenumi{\alph{enumi})}
\setcounter{enumi}{4}
\tightlist
\item
  Let w be a primitive element of \(\mathrm{E}(\lambda)\otimes\mathrm{E}(\lambda')\) . Assume that the weight \(\nu\) of w is distinct from \(\lambda+\lambda'\) . Show that there exists a chain \((\alpha_{1},\ldots,\alpha_{n})\) joining \(\lambda\) to \(\lambda'\) such that
\end{enumerate}

\[
\nu \leq \lambda + \lambda^ {\prime} - \alpha_ {1} - \dots - \alpha_ {n}.
\]

(Let C be the set of \(\alpha \in \mathrm{B}\) such that the coordinate of index \(\alpha\) of \(\lambda + \lambda' - \nu\) is \(\neq 0\) . Let D (resp. D') be the set of \(\alpha \in \mathrm{C}\) such that there exists \(\gamma_1, \ldots, \gamma_t \in \mathrm{C}\) satisfying \((\lambda|\gamma_1) \neq 0\) (resp. \((\lambda'|\gamma_1) \neq 0\) ), \((\gamma_1|\gamma_2) \neq 0, \ldots, (\gamma_{t-1}|\gamma_t) \neq 0\) , \((\gamma_t|\alpha) \neq 0\) . Let Y (resp. Y') be the set of weights of E( \(\lambda\) ) (resp. E( \(\lambda'\) )) of the form \(\lambda - \sum_{\alpha \in \mathrm{C}} k_\alpha\alpha\) (resp. \(\lambda' - \sum_{\alpha \in \mathrm{C}} k_\alpha\alpha\) ), with \(k_\alpha \in \mathbf{N}\) . Show that \(w\) belongs to \(\left(\sum_{\mu \in \mathrm{Y}}\mathrm{E}(\lambda)^{\mu}\right)\otimes \left(\sum_{\mu^{\prime}\in \mathrm{Y}^{\prime}}\mathrm{E}(\lambda^{\prime})^{\mu^{\prime}}\right)\) . Using \(a\) ) and the fact that \(\nu \neq \lambda +\lambda '\) , show that \(\mathrm{D}\cap \mathrm{D}'\neq \emptyset .\)

\(f\) ) Show the equivalence of the following properties:

\begin{enumerate}
\def\labelenumi{(\roman{enumi})}
\item
  \(\mathrm{E}(\lambda)\otimes \mathrm{E}(\lambda^{\prime})\) is isomorphic to \(\mathrm{E}(\lambda +\lambda^{\prime})\)
\item
  \(\mathrm{E}(\lambda) \otimes \mathrm{E}(\lambda')\) is a simple module.
\item
  There is no chain joining \(\lambda\) to \(\lambda'\) .
\item
  R is the direct sum of two root systems \(R_{1}\) and \(R_{1}^{\prime}\) such that \(\lambda \in \mathrm{P}(\mathrm{R}_{1})\) and \(\lambda^{\prime} \in \mathrm{P}(\mathrm{R}_{1}^{\prime})\) .
\item
  \(\mathfrak{g}\) is the product of two ideals \(\mathfrak{s}\) and \(\mathfrak{s}'\) such that \(\mathfrak{s}'\) .E(λ) = 0 and \(\mathfrak{s}\) .E(λ') = 0.
\end{enumerate}

(Use \(d\) ) to prove the equivalence of (ii) and (iii).) \(^{15}\)

\begin{enumerate}
\def\labelenumi{\arabic{enumi})}
\setcounter{enumi}{18}
\item
  We recall the notations of Prop. 10. Let \(\bar{k}\) be an algebraic closure of \(k\) . If \(x \in \mathfrak{g}\) , denote by \(\mathcal{X}_{\mathrm{E}}(x)\) (resp. \(\mathcal{X}_{\mathrm{F}}(x)\) , resp. \(\mathcal{X}_{\mathrm{G}}(x)\) ) the set of eigenvalues of \(x_{\mathrm{E}}\) (resp. \(x_{\mathrm{F}}\) , resp. \(x_{\mathrm{G}}\) ) in \(\bar{k}\) . Show that \(\mathcal{X}_{\mathrm{G}}(x) = \mathcal{X}_{\mathrm{E}}(x) + \mathcal{X}_{\mathrm{F}}(x)\) for all \(x \in \mathfrak{g}\) and that, when E and F are given, this property characterizes the simple \(\mathfrak{g}\) -module G up to isomorphism.
\item
  Assume that \(k = \mathbf{R}\) or \(\mathbf{C}\) . Let \(\Gamma\) be a simply-connected Lie group with Lie algebra \(\mathfrak{g}\) . Let \(\lambda, \mu, \mathrm{E}, \mathrm{F}, \mathrm{G}\) be as in Prop. 10. Let \((e_1, \ldots, e_n)\) (resp. \((f_1, \ldots, f_p)\) ) be a basis of \(\mathrm{E}\) (resp. \(\mathrm{F}\) ) consisting of eigenvectors of \(\mathfrak{h}\) , with \(e_1 \in \mathrm{E}^\lambda\) and \(f_1 \in \mathrm{F}^\mu\) . We can consider \(\mathrm{E}, \mathrm{F}, \mathrm{G}\) as \(\Gamma\) -modules. If \(\gamma \in \Gamma\) , denote by \(a_i(\gamma)\) the coordinate of \(\gamma.e_1\) with index \(i\) , and \(b_j(\gamma)\) the coordinate of \(\gamma.f_1\) with index \(j\) . Show that the function \(a_i b_j\) on \(\Gamma\) is not identically zero. Deduce that, for all \((i,j)\) , there exists an element of \(\mathrm{G} \subset \mathrm{E} \otimes \mathrm{F}\) whose coordinate of index \((i,j)\) is \(\neq 0\) ; for \(k = \mathbf{R}\) or \(\mathbf{C}\) , this gives a new proof of Prop. 10. Pass from this to the case \(k = \mathbf{Q}\) , and then to the case of an arbitrary field, cf.~Exerc. 4.
\item
  Let \(\lambda, \mu \in \mathrm{P}_{++}\) . Let E,F,G be simple \(\mathfrak{g}\) -modules of highest weights \(\lambda, \mu, \lambda + \mu\) , and let \(n\) be an integer \(\geq 1\) . If \(\omega\) is a weight of E of multiplicity \(n\) , then \(\omega + \mu\) is a weight of G of multiplicity \(\geq n\) . (We have \(\mathrm{G}^{\omega + \mu} \subset \bigoplus_{\nu + \sigma = \omega + \mu} \mathrm{E}^{\nu} \otimes \mathrm{F}^{\sigma}\) . If \(\dim \mathrm{G}^{\omega + \mu} < n\) , the projection of \(\mathrm{G}^{\omega + \mu}\) onto \(\mathrm{E}^{\omega} \otimes \mathrm{F}^{\mu}\) is of the form \(\mathrm{E}' \otimes \mathrm{F}^{\mu}\) , with \(\mathrm{E}'\) strictly contained in \(\mathrm{E}^{\omega}\) . Derive a contradiction from this by choosing adapted bases of E and F and imitating the proof of Prop. 10.)
\item
  Assume that \(\mathfrak{g}\) is simple, in other words that R is irreducible. Show that there exists a unique dominant weight \(\lambda \neq 0\) such that the set of weights of \(\mathrm{E}(\lambda)\) is \(\mathrm{W}.\lambda \cup \{0\}\) : we have \(\lambda = \alpha\) , where \(\alpha \in \mathrm{R}\) is such that \(H_{\alpha}\) is the highest root of \(\mathrm{R}^{\vee}\) . When all the roots are of the same length (cases \(\mathrm{A}_l,\mathrm{D}_l,\mathrm{E}_6,\mathrm{E}_7,\mathrm{E}_8\) ), we have \(\lambda = \tilde{\alpha}\) ; this is the only root that is a dominant weight; the corresponding representation is the adjoint representation of g. In the other cases, \(\lambda\) is the only root of minimum length that is a dominant weight; with the notations of Chap. VI, Plates, we have: \(\lambda = \varpi_{1}\) (type \(B_{l}\) ); \(\lambda = \varpi_{2}\) (type \(C_{l}\) ); \(\lambda = \varpi_{4}\) (type \(F_{4}\) ); \(\lambda = \varpi_{1}\) (type \(G_{2}\) ).
\item
  Let \(\mathrm{U}^0\) be the commutant of \(\mathfrak{h}\) in \(\mathrm{U}(\mathfrak{g})\) (cf.~§6, no. 4).
\end{enumerate}

\begin{enumerate}
\def\labelenumi{\alph{enumi})}
\item
  Let V be a (finite dimensional) g-module. Show that the \(V^{\lambda}\) , \(\lambda \in P\) , are stable under \(U^{0}\) , and that, if V is simple and \(V^{\lambda} \neq 0\) , \(V^{\lambda}\) is a simple \(U^{0}\) -module (use the decomposition \(\mathrm{U}(\mathfrak{g}) = \oplus\mathrm{U}^{\lambda}\) , loc. cit.).
\item
  Show that, for every element \(c \neq 0\) of \(U^{0}\) , there exists a finite dimensional simple representation \(\rho\) of \(U^{0}\) such that \(\rho(c) \neq 0\) . (Use a) and Chap. I, §7, Exerc. 3 a).
\end{enumerate}

\begin{enumerate}
\def\labelenumi{\arabic{enumi})}
\setcounter{enumi}{23}
\tightlist
\item
  Let \(A\) be a unital associative algebra and \(M\) the set of its finite codimensional two-sided ideals.
\end{enumerate}

\begin{enumerate}
\def\labelenumi{\alph{enumi})}
\tightlist
\item
  Let \(\mathrm{U}^*\) be the vector space dual of U. If \(\theta \in \mathrm{U}^*\) , show the equivalence of the following properties:
\end{enumerate}

\begin{enumerate}
\def\labelenumi{(\roman{enumi})}
\item
  there exists \(\mathfrak{m} \in \mathrm{M}\) such that \(\theta(\mathfrak{m}) = 0\) ;
\item
  there exist two finite families \((\theta_{i}^{\prime})\) and \((\theta_{i}^{\prime\prime})\) of elements of \(U^{*}\) such that
\end{enumerate}

\[
\theta (x y) = \sum_ {i} \theta_ {i} ^ {\prime} (x) \theta_ {i} ^ {\prime \prime} (y) \quad \text { for   all } x, y \in \mathrm{U}.
\]

The elements \(\theta\) with these properties form a subspace \(\mathrm{U}'\) of \(\mathrm{U}^*\) , which coincides with that denoted by \(\mathrm{B}'\) in Algebra, Chap. III, §11, Exerc. 27. There exists a unique coalgebra structure on \(\mathrm{U}'\) whose coproduct \(c: \mathrm{U}' \to \mathrm{U}' \otimes \mathrm{U}'\) is given by

\[
c (\theta) = \sum_ {i} \theta_ {i} ^ {\prime} \otimes \theta_ {i} ^ {\prime \prime},
\]

where \(\theta_i', \theta_i'' \in \mathrm{U}'\) are such that \(\theta(xy) = \sum_{i} \theta_i'(x) \theta_i''(y)\) for all \(x, y \in \mathrm{U}\) , cf.~(ii).

The coalgebra \(\mathrm{U}'\) is the union of the increasing filtration by the subspaces \((\mathrm{U} / \mathfrak{m})^*\) , \(\mathfrak{m} \in \mathrm{M}\) , which are finite dimensional. The dual of \(\mathrm{U}'\) can be identified with the algebra \(\hat{\mathrm{U}} = \varprojlim \mathrm{U} / \mathfrak{m}\) ; if \(\hat{\mathrm{U}}\) is given the projective limit of the discrete topologies on the \(\mathrm{U} / \mathfrak{m}, \mathfrak{m} \in \mathrm{M}\) , the continuous linear forms on \(\hat{\mathrm{U}}\) are given by the elements of \(\mathrm{U}'\) .


\begin{enumerate}
\def\labelenumi{\alph{enumi})}
\setcounter{enumi}{1}
\tightlist
\item
  Let E be a finite dimensional left U-module. Its annihilator \(m_{E}\) belongs to M; the composite \(\hat{U} \to U/m_{E} \to \text{End}(E)\) gives E the structure of a left \(\hat{U}\) -module. If F is a finite dimensional left U-module, a linear map \(f : E \to F\) is a U-homomorphism if and only if it is a \(\hat{U}\) -homomorphism. If \(a \in E, b \in E^{*}\) , the linear form \(\theta_{a,b} : x \mapsto \langle xa, b \rangle\) belongs to \(U'\) , and
\end{enumerate}

\[
\langle x, \theta_ {a, b} \rangle = \langle x a, b \rangle \quad \text { for   all } x \in \hat {\mathrm{U}}.
\]

The \(\theta_{a,b}\) (for varying E, \(a,b\) ) generate the \(k\) -vector space U'.

\begin{enumerate}
\def\labelenumi{\alph{enumi})}
\setcounter{enumi}{2}
\tightlist
\item
  Let \(X_{U}\) be the set of isomorphism classes of finite dimensional left U-modules. For all \(E \in X_{U}\) , let \(u_{E}\) be a k-linear endomorphism of E; assume that
\end{enumerate}

\[
f \circ u _ {\mathrm{E}} = u _ {\mathrm{F}} \circ f \quad \text { for   all   E, F } \in X_{U} \text { and } f \in \operatorname{Hom} _ {\mathrm{U}} (\mathrm{E,F}).
\]

Show that there exists a unique element \(x \in \hat{U}\) such that \(x_{E} = u_{E}\) for all \(E \in X_{U}\) . (Reduce to the case in which U is finite dimensional.)

¶ 25) Let a be a Lie algebra and U its enveloping algebra. Apply the definitions and results of Exerc. 24 to the algebra U. In particular, \(U' \subset U^*\) and the dual of \(U'\) can be identified with the algebra \(\hat{U} = \varprojlim U/m\) ; the canonical map \(U \to \hat{U}\) is injective (Chap. I, §7, Exerc. 3).

\begin{enumerate}
\def\labelenumi{\alph{enumi})}
\tightlist
\item
  The coalgebra structure of U (Chap. II, §1, no. 4) defines an algebra structure on the dual U* (cf.~Chap. II, §1, no. 5, Prop. 10, and Algebra, Chap. III, §11, no. 2). If E, F are finite dimensional a-modules,
\end{enumerate}

\[
\theta_ {a, b} \circ \theta_ {c, d} = \theta_ {a \otimes c, b \otimes d} \text { for } a \in \mathrm{E}, b \in \mathrm{E} ^ {*}, c \in \mathrm{F}, d \in \mathrm{F} ^ {*}.
\]

Deduce that \(U'\) is a subalgebra of \(U^{*}\) . The coalgebra and algebra structures of \(U'\) make it a commutative bigebra (Algebra, Chap. III, §11, no. 4).

\begin{enumerate}
\def\labelenumi{\alph{enumi})}
\setcounter{enumi}{1}
\tightlist
\item
  Let \(x\) be an element of \(\hat{\mathrm{U}}\) ; identify \(x\) with a linear form \(\mathrm{U}' \to k\) . Prove the equivalence of the following properties:
\end{enumerate}

\begin{enumerate}
\def\labelenumi{(\roman{enumi})}
\item
  \(x\) is a homomorphism of algebras from \(\mathrm{U}'\) to \(k\) .
\item
  \(x_{E \otimes F} = x_{E} \otimes x_{F}\) for all finite dimensional a-modules E and F.
\end{enumerate}

(Prove first that (ii) is equivalent to

\[
\left(\mathrm{ii} ^ {\prime}\right) x \left(\theta_ {a \otimes c, b \otimes d}\right) = x \left(\theta_ {a, b}\right) x \left(\theta_ {c, d}\right) \text {if} a \in \mathrm{E}, b \in \mathrm{E} ^ {*}, c \in \mathrm{F}, d \in \mathrm{F} ^ {*},
\]

and use the fact that the \(\theta_{a,b}\) generate \(\mathrm{U}'\) .)

\(c)\) Let \(x\) be an element of \(\hat{\mathbf{U}}\) satisfying conditions (i) and (ii) of \(b\) ). Show the equivalence of the following conditions:

\begin{enumerate}
\def\labelenumi{(\roman{enumi})}
\setcounter{enumi}{2}
\item
  \(x\) takes the unit element of \(\mathrm{U}'\) to the unit element of \(k\) .
\item
  \(x \neq 0.\)
\item
  If \(k\) is given the trivial \(\mathfrak{a}\) -module structure, we have \(x_{k} = \mathrm{Id}\) .
\item
  \(x_{\mathrm{E}}\) is invertible for all E.
\end{enumerate}

(The equivalence (iii) \(\Longleftrightarrow\) (iv) follows from the fact that \(x\) is a homomorphism of algebras. On the other hand, \(x_{k} = \lambda \mathrm{Id}\) , with \(\lambda \in k\) . Using the \(\mathfrak{a}\) -isomorphism \(k \otimes \mathrm{E} \to \mathrm{E}\) , deduce that \(\lambda x_{\mathrm{E}} = x_{\mathrm{E}}\) for all E, and, in particular, that \(\lambda^2 = \lambda\) by taking \(\mathrm{E} = k\) . The case \(\lambda = 1\) corresponds to \(x \neq 0\) , hence (iv) \(\Longleftrightarrow\) (v), and (vi) \(\Longrightarrow\) (v). To prove that (v) \(\Longrightarrow\) (vi), show that, if F is the dual of E, then \(^t x_{\mathrm{F}} \circ x_{\mathrm{E}} = \lambda \mathrm{Id}_{\mathrm{E}}\) .)

\begin{enumerate}
\def\labelenumi{\alph{enumi})}
\setcounter{enumi}{3}
\tightlist
\item
  Let G be the set of elements of \(\hat{U}\) satisfying conditions (i) to (vi) above. Show that G is a subgroup of the group of invertible elements of \(\hat{U}\) .
\end{enumerate}

Let \(x \in G\) . If E is a finite dimensional \(\mathfrak{a}\) -module, then \(x_{\mathrm{E}} \in \mathbf{GL}(\mathrm{E})\) . This applies in particular to \(\mathrm{E} = \mathfrak{a}\) , equipped with the adjoint representation; this gives an element \(x_{\mathfrak{a}} \in \mathbf{GL}(\mathfrak{a})\) . Show that \(x_{\mathfrak{a}}\) is an automorphism of \(\mathfrak{a}\) (use the \(\mathfrak{a}\) -homomorphism \(\mathfrak{a} \otimes \mathfrak{a} \to \mathfrak{a}\) given by the bracket). This gives a homomorphism \(v: \mathrm{G} \to \operatorname{Aut}(\mathfrak{a})\) . If \(\mathrm{E}\) is a finite dimensional \(\mathfrak{a}\) -module,

\[
x _ {\mathrm{E}} (y. e) = v (x) (y). x _ {\mathrm{E}} (e) \quad \mathrm{if} y \in \mathfrak {a}, e \in \mathrm{E}
\]

(use the \(\mathfrak{a}\) -homomorphism \(\mathfrak{a} \otimes \mathrm{E} \to \mathrm{E}\) given by the operation of \(\mathfrak{a}\) on \(\mathrm{E}\) ).

\begin{enumerate}
\def\labelenumi{\alph{enumi})}
\setcounter{enumi}{4}
\tightlist
\item
  The principal anti-automorphism \(\sigma\) of U extends by continuity to \(\hat{\mathbf{U}}\) . Its transpose leaves \(\mathrm{U}'\) stable and induces an inversion on \(\mathrm{U}'\) (Algebra, Chap. III, §11, Exerc. 4). If \(x \in \mathrm{G}\) , then \(\sigma(x) = x^{-1}\) .
\end{enumerate}

\(f\) ) Assume that \(\mathfrak{a}\) is semi-simple \(^{16}\) . Let \(n\) be a nilpotent element of \(\mathfrak{a}\) . Then there exists a unique element \(e^n\) of \(\mathrm{G}\) such that \((e^n)_{\mathrm{E}} = \exp (n_{\mathrm{E}})\) for every finite dimensional \(\mathfrak{a}\) -module E. We have \(v(e^n) = \exp (\operatorname {ad}n)\in \operatorname {Aut}(\mathfrak{a})\) , so \(\operatorname {Aut}_e(\mathfrak{a})\subset v(\mathrm{G})\) .

Show that, if b is a subalgebra of a consisting of nilpotent elements, then

\[
e ^ {n}. e ^ {m} = e ^ {\mathrm{H} (n, m)} \quad \text { for } n, m \in \mathfrak {b},
\]

where H denotes the Hausdorff series (Chap. II, §6).

¶ 26) Apply the notations and results of Exerc. 25 to the case in which a = g (split case).

\begin{enumerate}
\def\labelenumi{\alph{enumi})}
\item
  Let \(x \in G\) and let \(\sigma = v(x)\) be its image in \(\operatorname{Aut}(\mathfrak{g})\) . If \(\rho\) is a representation of g, \(\rho\) and \(\rho \circ \sigma\) are equivalent. Deduce (cf.~no. 2, Remark 1) that \(\sigma\) belongs to \(\operatorname{Aut}_{0}(\mathfrak{g})\) . Extend this result to arbitrary semi-simple algebras.
\item
  Let \(\varphi \in \mathrm{T}_{\mathrm{P}} = \mathrm{Hom}(\mathrm{P}, k^{*})\) , where \(\mathrm{P} = \mathrm{P}(\mathrm{R})\) . If \(\mathrm{E}\) is a \(\mathfrak{g}\) -module, let \(\varphi_{\mathrm{E}}\) be the endomorphism of \(\mathrm{E}\) whose restriction to each \(\mathrm{E}^{\lambda}\) ( \(\lambda \in \mathrm{P}\) ) is the homothety of ratio \(\varphi(\lambda)\) . Show that there exists a unique element \(t(\varphi) \in \mathrm{G}\) such that \(t(\varphi)_{\mathrm{E}} = \varphi_{\mathrm{E}}\) for all \(\mathrm{E}\) (Use Exerc. 24 c), and the characterizations (ii) and (vi) of Exerc. 25.) This gives a homomorphism \(t: \mathrm{T}_{\mathrm{P}} \to \mathrm{G}\) . Show that \(t\) is injective. We use this to identify \(\mathrm{T}_{\mathrm{P}}\) with a subgroup of \(\mathrm{G}\) . The composite \(\mathrm{T}_{\mathrm{P}} \to \mathrm{G} \to \mathrm{Aut}(\mathfrak{g})\) is the homomorphism denoted by \(f \circ q\) in §5, no. 2.
\item
  Let \(x \in G\) be such that \(\sigma = v(x)\) belongs to the subgroup \(f(\mathrm{T}_{\mathrm{Q}})\) of \(\mathrm{Aut}_{0}(\mathfrak{g})\) (§5, no. 2), in other words, such that it operates trivially on h; denote the element of \(T_{Q} = \mathrm{Hom}(Q, k^{*})\) corresponding to x by \(\psi\) . We are going to show that x belongs to \(T_{P}\) . Prove successively:
\end{enumerate}

\(c_{1}\) ) If E is a \(\mathfrak{g}\) -module, \(x_{\mathrm{E}}\) is an \(\mathfrak{h}\) -endomorphism of E.

(Use the g-homomorphism \(g \otimes E \to E\) , and the fact that x operates trivially on h.) In particular, the \(E^{\mu}\) are stable under \(x_{E}\) .

\(c_{2}\) ) There exists \(\varphi \in \mathrm{T}_{\mathrm{P}}\) such that, for every \(\mathfrak{g}\) -module E, and every primitive element \(e\) of E of weight \(\lambda\) , we have \(x_{\mathrm{E}}e = \varphi (\lambda)e\) .

(Choose \(\varphi\) such that this relation is true when E is a fundamental module \(\mathrm{E}(\varpi_{\alpha})\) . Deduce the case of the \(\mathrm{E}(\lambda)\) , \(\lambda \in \mathrm{P}_{++}\) , by using the embedding of such a module in a tensor product of the \(\mathrm{E}(\varpi_{\alpha})\) . Pass from this to the general case.)

\(c_{3}\) ) Choose \(\varphi\) as in \(c_{2}\) ). Let E be a simple g-module of highest weight \(\lambda\) , and let \(\mu\) be a weight of E; then \(\lambda - \mu \in Q\) . Show that the restriction of \(x_{E}\) to \(E^{\mu}\) is the homothety of ratio \(\varphi(\lambda)\psi(\mu - \lambda)\) . (Same method as for \(c_{1}\) .)

\(c_{4}\) ) If \(\lambda, \mu \in P_{++}\) , and if \(\alpha \in B\) is not orthogonal to either \(\lambda\) or \(\mu\) , then

\[
\varphi (\lambda + \mu - \alpha) = \varphi (\lambda + \mu) \psi (- \alpha),
\]

so \(\varphi (\alpha) = \psi (\alpha)\) . (Use \(c_{2}\) ), \(c_{3}\) ), and the embedding of \(\operatorname {E}(\lambda +\mu -\alpha)\) in \(\operatorname {E}(\lambda)\otimes \operatorname {E}(\mu)\) , cf.~Exerc. 17.)

\(c_{5})\) Deduce from \(c_{4})\) that \(\varphi|\mathbf{Q}=\psi\) , and use \(c_{3})\) to deduce that \(x=t(\varphi)\) . \(d)\) Identify \(\mathrm{T}_{\mathrm{Q}}\) with a subgroup of \(\mathrm{Aut}_{0}(\mathfrak{g})\) by means of \(f\) . By \(a)\) , we have

\[
\operatorname{Aut} _ {e} (\mathfrak {g}) \subset v (\mathrm{G}) \subset \operatorname{Aut} _ {0} (\mathfrak {g})
\]

and, by \(c\) ), \(v(\mathrm{G}) \cap \mathrm{T}_{\mathrm{Q}} = \mathrm{Im}(\mathrm{T}_{\mathrm{P}})\) . Deduce (cf.~§5, no. 3) that \(\mathrm{Aut}_e(\mathfrak{g}) \cap \mathrm{T}_{\mathrm{Q}} = \mathrm{Im}(\mathrm{T}_{\mathrm{P}})\) and that \(f(\mathrm{G}) = \mathrm{Aut}_e(\mathfrak{g})\) . The canonical map

\[
\iota : \mathrm{T} _ {\mathrm{Q}} / \operatorname{Im} (\mathrm{T} _ {\mathrm{P}}) \to \operatorname{Aut} _ {0} (\mathfrak {g}) / \operatorname{Aut} _ {e} (\mathfrak {g})
\]

is therefore an isomorphism.

\begin{enumerate}
\def\labelenumi{\alph{enumi})}
\setcounter{enumi}{4}
\tightlist
\item
  The kernel of \(v: \mathrm{G} \to \mathrm{Aut}_e(\mathfrak{g})\) is equal to the kernel of \(\mathrm{T}_{\mathrm{P}} \to \mathrm{T}_{\mathrm{Q}}\) ; it is isomorphic to
\end{enumerate}

\[
\operatorname{Hom} (\mathrm{P} / \mathrm{Q}, k ^ {*});
\]

this is a finite abelian group contained in the centre of G, and its order divides \((P:Q)\) ; if k is algebraically closed, it is isomorphic to the dual of P/Q (Algebra, Chap. VII, §4, no. 8).

\(f)\) Let \(\alpha \in \mathrm{R}\) , \(X_{\alpha} \in \mathfrak{g}^{\alpha}\) and \(X_{-\alpha} \in \mathfrak{g}^{-\alpha}\) be such that \([X_{\alpha}, X_{-\alpha}] = -H_{\alpha}\) , and let \(\rho_{\alpha}\) be the corresponding representation of \(\mathfrak{sl}(2, k)\) on \(\mathfrak{g}\) . If E is a \(\mathfrak{g}\) -module, deduce (§1, no. 4) a representation of \(\mathbf{SL}(2, k)\) on E, and hence (Exerc. 25 b), c)) a homomorphism

\[
\varphi_ {\alpha}: \mathbf {SL} (2, k) \to \mathrm{G}.
\]

Show that \(\operatorname{Im}(\varphi_{\alpha})\) contains the elements of \(T_{P}\) of the form \(\lambda \mapsto t^{\lambda(H_{\alpha})}, t \in k^{*}\) . Deduce that the \(\operatorname{Im}(\varphi_{\alpha}), \alpha \in B\) , generate G (show first that the group they generate contains \(T_{P}\) ). In particular, G is generated by the \(e^{n}\) , with \(n \in g^{\alpha}, \alpha \in B \cup -B\) . The derived group of G is equal to G.

\(g\) ) If a subgroup \(\mathbf{G}'\) of \(\mathbf{G}\) is such that \(v(\mathrm{G}') = \mathrm{Aut}_e(\mathfrak{g})\) , then \(\mathrm{G}' = \mathrm{G}\) (use \(f\) ). \(h\) ) Let \(\mathrm{E}\) be a faithful \(\mathfrak{g}\) -module, and let \(\Gamma\) be the subgroup of \(\mathrm{P}\) generated by the weights of \(\mathrm{E}\) . Then \(\mathrm{P} \supset \Gamma \supset \mathrm{Q}\) , cf.~Exerc. 5. Show that the kernel of the canonical homomorphism \(\mathrm{G} \to \mathbf{GL}(\mathrm{E})\) is equal to the subgroup of \(T_{P}\) consisting of the elements \(\varphi\) whose restriction to \(\varGamma\) is trivial. In particular, if \(\varGamma = P\) the homomorphism \(\mathrm{G} \to \mathbf{GL}(\mathrm{E})\) is injective. If \(\varGamma = Q\) , this homomorphism factorizes as \(\mathrm{G} \stackrel{v}{\longrightarrow} \operatorname{Aut}_{e}(\mathfrak{g}) \longrightarrow \mathbf{GL}(\mathrm{E})\) , and the homomorphism \(\operatorname{Aut}_{e}(\mathfrak{g}) \to \mathbf{GL}(\mathrm{E})\) is injective.

\begin{enumerate}
\def\labelenumi{\arabic{enumi})}
\setcounter{enumi}{26}
\tightlist
\item
  Let \(\Omega = \mathrm{P} / \mathrm{Q}\) . If \(\omega \in \Omega\) , and if \(\mathrm{E}\) is a \(\mathfrak{g}\) -module, denote by \(\mathrm{E}_{\omega}\) the direct sum of the \(\mathrm{E}^{\lambda}\) , for \(\lambda \in \omega\) . We have \(\mathrm{E} = \bigoplus_{\omega \in \Omega} \mathrm{E}_{\omega}\) .
\end{enumerate}

\begin{enumerate}
\def\labelenumi{\alph{enumi})}
\tightlist
\item
  Show that \(E_{\omega}\) is a g-submodule of E. We have \((\mathrm{E}^{*})_{\omega} = (\mathrm{E}_{-\omega})^{*}\) and
\end{enumerate}

\[
(\mathrm{E} \otimes \mathrm{F}) _ {\omega} = \bigoplus_ {\alpha + \beta = \omega} \mathrm{E} _ {\alpha} \otimes \mathrm{F} _ {\beta}
\]

if F is another g-module.

\begin{enumerate}
\def\labelenumi{\alph{enumi})}
\setcounter{enumi}{1}
\item
  Let \(\chi\in\mathrm{Hom}(\varOmega,k^{*})=\mathrm{Ker}(\mathrm{T}_{\mathrm{P}}\to\mathrm{T}_{\mathrm{Q}})\) . Identify \(\chi\) with an element of the kernel of \(f:\mathrm{G}\to\mathrm{Aut}_{e}(\mathfrak{g})\) , cf.~Exerc. 26 e). Show that the operation of \(\chi\) on \(E_{\omega}\) is the homothety of ratio \(\chi(\omega)\) .
\item
  What are the \(\mathrm{E}_{\omega}\) when \(\mathfrak{g} = \mathfrak{sl}(2,k)\) ?
\end{enumerate}

